\hypersetup{pageanchor=false}
\begin{titlepage}
\centering
\vspace*{8mm}
\includegraphics[width=105mm]{figures/iscte-university-logo.png}\par
\vspace{28mm}
\parbox{140mm}{\centering\LARGE\bfseries\thesistitle}\par
\vspace{20mm}
{\large\thesisauthor\par}
\vspace{14mm}
{\large\thesisdegree\par}
\vspace{14mm}
Supervisors:\par
\vspace{3mm}
\begin{minipage}{0.96\textwidth}
\centering
\thesissupervisors
\end{minipage}\par
\vfill
{\thesisdate\par}

\end{titlepage}
\cleardoublepage
\begin{titlepage}
\centering
\vspace*{0mm}
\includegraphics[width=72mm]{figures/iscte-department-logo.png}\par
\vspace{18mm}
\parbox{140mm}{\centering\LARGE\bfseries\thesistitle}\par
\vspace{13mm}
{\large\thesisauthor\par}
\vspace{10mm}
{\thesisdegree\par}
\vspace{14mm}
Supervisors:\par
\vspace{3mm}
\begin{minipage}{0.96\textwidth}
\centering
\thesissupervisors
\end{minipage}\par
\vfill
\thesisinstitution\par
\thesisdate\par
\end{titlepage}
\cleardoublepage
\pagenumbering{roman}
\hypersetup{pageanchor=true}
\chapter*{Acknowledgements}
\addcontentsline{toc}{chapter}{Acknowledgements}
\draftfield{Acknowledgements}

\chapter*{Resumo}
\addcontentsline{toc}{chapter}{Resumo}
\begin{otherlanguage}{portuguese}
Durante a resposta a um incidente de cibersegurança, encontrar informação não basta: é necessário reconhecer a orientação aplicável, compreender as suas condições e decidir como a utilizar. Quando procedimentos, relatórios de incidentes e informação sobre ameaças se encontram dispersos, a experiência acumulada pode ser difícil de recuperar no momento em que é necessária. Neste trabalho, investigamos como um sistema de gestão de conhecimento pode utilizar \emph{machine learning} para apoiar essa consulta e tornar o conhecimento disponível mais útil para a decisão.

Adotamos uma abordagem de investigação orientada à conceção, ligando a análise do problema à definição de requisitos, ao desenvolvimento de uma solução e à sua avaliação. Consideramos o ciclo de gestão do conhecimento, desde a recolha e organização até à atualização e reutilização. O papel da aprendizagem automática é estudado na recuperação de informação relevante e na formulação de orientações baseadas em fontes. A responsabilidade pela decisão mantém-se com o analista.

Definimos a otimização da resposta a incidentes através da qualidade da informação recuperada, da adequação das orientações e do esforço necessário para realizar as tarefas. A avaliação relaciona estas dimensões com a experiência de utilização, sem reduzir o desempenho a uma única medida de rapidez. Procuramos, assim, aproximar a gestão de conhecimento das necessidades da resposta a incidentes, preservando a possibilidade de consultar as fontes e respeitando a confidencialidade da informação.
% Finalise the empirical findings sentence after importing the researcher's recorded measurements.

\medskip
\noindent\textbf{Palavras-chave:} gestão de conhecimento; aprendizagem automática; resposta a incidentes; cibersegurança; recuperação de informação.
\end{otherlanguage}

\chapter*{Abstract}
\addcontentsline{toc}{chapter}{Abstract}
During a cybersecurity incident, finding information is only part of the task. An analyst must identify applicable guidance, understand its conditions and decide how to use it. When procedures, incident reports and threat information are dispersed, accumulated experience may be difficult to recover when it is needed. We investigate how a knowledge-management system can use machine learning to support this process and make documented knowledge more useful for incident-response decisions.

We adopt a design-science approach that connects the problem analysis with requirements, development and evaluation. Our scope covers the knowledge-management cycle, including collection, organisation, maintenance and reuse. We examine machine learning as a means of retrieving relevant information and formulating source-based guidance, while leaving responsibility for operational decisions with the analyst.

We define optimisation through the quality of retrieved information, the suitability of the guidance and the effort required to complete a task. The evaluation connects these dimensions with the experience of using the system, rather than reducing performance to a single speed measure. This perspective places knowledge management within the practical demands of incident response: making relevant experience accessible, retaining the means to examine its sources and respecting information confidentiality.
% Finalise the empirical findings sentence after importing the researcher's recorded measurements.

\medskip
\noindent\textbf{Keywords:} knowledge management; machine learning; incident response; cybersecurity; information retrieval.
